% GENERATED FILE. Edit canonical lessons, then run npm run generate:books.
% canonical-source-hash: fnv1a64:3ac21af6ffdb51d0
% canonical-lessons: UR-C478-chubhna, UR-C478-chira, UR-C478-tanka, UR-C478-pip, UR-C478-kuredna

\chapter{Cuts and Stitches}
\label{ch:ur-cuts-and-stitches}
\hlchaptermodality{478}

\begin{chapteropening}
\textbf{By the end of this chapter:} \emph{I can say how someone got a small injury.}

Complete the last lesson of chapter 478: I can say how someone got a small injury.
\end{chapteropening}

\section[\texorpdfstring{\ur{چبھنا} (chubhnā)}{چبھنا (chubhnā)}]{\ur{چبھنا} (chubhnā) — to prick; to stick in}

\label{lesson:UR-C478-chubhna}



\begin{quote}
Before the new one: say the Urdu for a paper packet, then the Urdu for to toss into the mouth.
\end{quote}

\subsection*{You'll want to know: \ur{چبھنا}}
\textbf{\ur{چبھنا}} — \emph{chubhnā} — \textquotedblleft{}to prick; to stick in\textquotedblright{}.

\begin{quote}
\ur{میری} \ur{انگلی} \ur{میں} \ur{کانٹا} \ur{چبھ} \ur{گیا۔}
\end{quote}

\begin{quote}
\emph{merī unglī meṅ kā̃ṭā chubh gayā.}
\end{quote}

\begin{quote}
A thorn stuck in my finger.
\end{quote}

Past (it …): \ur{وہ} \ur{کانٹا} \ur{چبھا}, \emph{voh kā̃ṭā chubhā}.

\begin{cousinweb}
\textbf{\ur{چبھنا}} (\emph{chubhnā}) is inherited from \emph{cubhyatē}, \textquotedblleft{}is pierced\textquotedblright{}, a form CDIAL reconstructs (CDIAL 4867).
\end{cousinweb}

\begin{grammarlens}[title={What you've built}]
One more word for describing what you lived and felt.
\end{grammarlens}

\subsection*{Your turn}
\begin{itemize}
\raggedright
  \item \emph{Say it:} \emph{chubhnā}
  \item \emph{Say it:} \emph{chubhnā}, once more
  \item \emph{Say it:} \emph{merī unglī meṅ kā̃ṭā chubh gayā.} — A thorn stuck in my finger.
  \item \emph{Recall:} say the Urdu for a paper packet, then the Urdu for to toss into the mouth, then say \emph{chubhnā} again
\end{itemize}

\begin{tcolorbox}[breakable,colback=teal!4,colframe=teal!35!black,title={Before you move on}]
What does \ur{چبھنا} mean? (\textquotedblleft{}To prick; to stick in\textquotedblright{}.) Say it once more.
\end{tcolorbox}
\section[\texorpdfstring{\ur{چیرا} (chīrā)}{چیرا (chīrā)}]{\ur{چیرا} (chīrā) — an incision; a slit}

\label{lesson:UR-C478-chira}



\begin{quote}
Before the new one: say the Urdu for to toss into the mouth, then the Urdu for to prick.
\end{quote}

\subsection*{You'll want to know: \ur{چیرا}}
\textbf{\ur{چیرا}} — \emph{chīrā} — \textquotedblleft{}an incision; a slit\textquotedblright{}.

\begin{quote}
\ur{ڈاکٹر} \ur{نے} \ur{ہاتھ} \ur{پر} \ur{ایک} \ur{چھوٹا} \ur{چیرا} \ur{دیا۔}
\end{quote}

\begin{quote}
\emph{ḍākṭar ne hāth par ek chhoṭā chīrā diyā.}
\end{quote}

\begin{quote}
The doctor made a small incision in the hand.
\end{quote}

Masculine.

\begin{cousinweb}
\textbf{\ur{چیرا}} (\emph{chīrā}) is inherited: CDIAL lists the Hindi cousin \emph{cīrā}, \textquotedblleft{}a cut, a slit\textquotedblright{}, under Sanskrit \emph{cīra}, \textquotedblleft{}a strip of bark or cloth\textquotedblright{} (CDIAL 4843).
\end{cousinweb}

\begin{grammarlens}[title={What you've built}]
One more word for describing what you lived and felt.
\end{grammarlens}

\subsection*{Your turn}
\begin{itemize}
\raggedright
  \item \emph{Say it:} \emph{chīrā}
  \item \emph{Say it:} \emph{chīrā}, once more
  \item \emph{Say it:} \emph{ḍākṭar ne hāth par ek chhoṭā chīrā diyā.} — The doctor made a small incision in the hand.
  \item \emph{Recall:} say the Urdu for to toss into the mouth, then the Urdu for to prick, then say \emph{chīrā} again
\end{itemize}

\begin{tcolorbox}[breakable,colback=teal!4,colframe=teal!35!black,title={Before you move on}]
What does \ur{چیرا} mean? (\textquotedblleft{}An incision; a slit\textquotedblright{}.) Say it once more.
\end{tcolorbox}
\section[\texorpdfstring{\ur{ٹانکا} (ṭāṅkā)}{ٹانکا (ṭāṅkā)}]{\ur{ٹانکا} (ṭāṅkā) — a stitch}

\label{lesson:UR-C478-tanka}



\begin{quote}
Before the new one: say the Urdu for to prick, then the Urdu for an incision.
\end{quote}

\subsection*{You'll want to know: \ur{ٹانکا}}
\textbf{\ur{ٹانکا}} — \emph{ṭāṅkā} — \textquotedblleft{}a stitch\textquotedblright{}.

\begin{quote}
\ur{زخم} \ur{پر} \ur{تین} \ur{ٹانکے} \ur{لگے۔}
\end{quote}

\begin{quote}
\emph{zakhm par tīn ṭāṅke lage.}
\end{quote}

\begin{quote}
The wound needed three stitches.
\end{quote}

Masculine.

\begin{cousinweb}
\textbf{\ur{ٹانکا}} (\emph{ṭāṅkā}) is inherited: CDIAL lists the Hindi cousin \emph{ṭā̃kā}, \textquotedblleft{}a stitch, a join\textquotedblright{}, under Sanskrit \emph{ṭáṅkati}, \textquotedblleft{}ties\textquotedblright{} (CDIAL 5432).
\end{cousinweb}

\begin{grammarlens}[title={What you've built}]
One more word for describing what you lived and felt.
\end{grammarlens}

\subsection*{Your turn}
\begin{itemize}
\raggedright
  \item \emph{Say it:} \emph{ṭāṅkā}
  \item \emph{Say it:} \emph{ṭāṅkā}, once more
  \item \emph{Say it:} \emph{zakhm par tīn ṭāṅke lage.} — The wound needed three stitches.
  \item \emph{Recall:} say the Urdu for to prick, then the Urdu for an incision, then say \emph{ṭāṅkā} again
\end{itemize}

\begin{tcolorbox}[breakable,colback=teal!4,colframe=teal!35!black,title={Before you move on}]
What does \ur{ٹانکا} mean? (\textquotedblleft{}A stitch\textquotedblright{}.) Say it once more.
\end{tcolorbox}
\section[\texorpdfstring{\ur{پیپ} (pīp)}{پیپ (pīp)}]{\ur{پیپ} (pīp) — pus}

\label{lesson:UR-C478-pip}



\begin{quote}
Before the new one: say the Urdu for an incision, then the Urdu for a stitch.
\end{quote}

\subsection*{You'll want to know: \ur{پیپ}}
\textbf{\ur{پیپ}} — \emph{pīp} — \textquotedblleft{}pus\textquotedblright{}.

\begin{quote}
\ur{زخم} \ur{میں} \ur{پیپ} \ur{ہے۔}
\end{quote}

\begin{quote}
\emph{zakhm meṅ pīp hai.}
\end{quote}

\begin{quote}
There is pus in the wound.
\end{quote}

Feminine.

\begin{cousinweb}
\textbf{\ur{پیپ}} (\emph{pīp}) is inherited: CDIAL lists the Hindi cousin \emph{pīp}, \textquotedblleft{}pus\textquotedblright{}, under Sanskrit \emph{pīvan}, \textquotedblleft{}well-nourished, swelling\textquotedblright{} (CDIAL 8242).
\end{cousinweb}

\begin{grammarlens}[title={What you've built}]
One more word for describing what you lived and felt.
\end{grammarlens}

\subsection*{Your turn}
\begin{itemize}
\raggedright
  \item \emph{Say it:} \emph{pīp}
  \item \emph{Say it:} \emph{pīp}, once more
  \item \emph{Say it:} \emph{zakhm meṅ pīp hai.} — There is pus in the wound.
  \item \emph{Recall:} say the Urdu for an incision, then the Urdu for a stitch, then say \emph{pīp} again
\end{itemize}

\begin{tcolorbox}[breakable,colback=teal!4,colframe=teal!35!black,title={Before you move on}]
What does \ur{پیپ} mean? (\textquotedblleft{}Pus\textquotedblright{}.) Say it once more.
\end{tcolorbox}
\section[\texorpdfstring{\ur{کریدنا} (kurednā)}{کریدنا (kurednā)}]{\ur{کریدنا} (kurednā) — to pick at; to scrape at}

\label{lesson:UR-C478-kuredna}



\begin{quote}
Before the new one: say the Urdu for a stitch, then the Urdu for pus.
\end{quote}

\subsection*{You'll want to know: \ur{کریدنا}}
\textbf{\ur{کریدنا}} — \emph{kurednā} — \textquotedblleft{}to pick at; to scrape at\textquotedblright{}.

\begin{quote}
\ur{زخم} \ur{کو} \ur{نہیں} \ur{کریدنا} \ur{چاہیے۔}
\end{quote}

\begin{quote}
\emph{zakhm ko nahīṅ kurednā chāhiye.}
\end{quote}

\begin{quote}
You should not pick at a wound.
\end{quote}

Past, with \ur{نے} (he or she …): \ur{اس} \ur{نے} \ur{زخم} \ur{کو} \ur{کریدا}, \emph{us ne zakhm ko kuredā}.

\begin{cousinweb}
\textbf{\ur{کریدنا}} (\emph{kurednā}) is inherited: CDIAL lists the Hindi cousin \emph{kurednā}, \textquotedblleft{}to scrape\textquotedblright{}, under \emph{kura}, \textquotedblleft{}scraping\textquotedblright{}, a form it reconstructs (CDIAL 3319).
\end{cousinweb}

\begin{grammarlens}[title={What you've built}]
One more word for describing what you lived and felt.
\end{grammarlens}

\subsection*{Your turn}
\begin{itemize}
\raggedright
  \item \emph{Say it:} \emph{kurednā}
  \item \emph{Say it:} \emph{kurednā}, once more
  \item \emph{Say it:} \emph{zakhm ko nahīṅ kurednā chāhiye.} — You should not pick at a wound.
  \item \emph{Recall:} say the Urdu for a stitch, then the Urdu for pus, then say \emph{kurednā} again
\end{itemize}

\begin{tcolorbox}[breakable,colback=teal!4,colframe=teal!35!black,title={Before you move on}]
What does \ur{کریدنا} mean? (\textquotedblleft{}To pick at; to scrape at\textquotedblright{}.) Say it once more.
\end{tcolorbox}
